\documentclass[11pt]{article}
\usepackage[utf8]{inputenc}
\usepackage[T1]{fontenc}
\usepackage[spanish]{babel}
\usepackage[a4paper,margin=2.5cm]{geometry}
\usepackage{amsmath,amssymb}
\usepackage{siunitx}

\title{Acoplamiento crítico $K_c$ --- Oscilador de Kuramoto}
\author{}
\date{}

\begin{document}
\maketitle
\vspace{-3em}

\section*{Modelo}
\[
\frac{d\theta_i}{dt} = \omega_i + K\sum_j A_{ij}\,\sin(\theta_j-\theta_i),
\qquad \omega_i\sim\mathcal{N}(\mu,\sigma^2).
\]
Parámetro de orden $r(t)=\bigl|\tfrac{1}{N}\sum_j e^{\,i\theta_j}\bigr|$, con valor
estacionario $r_\infty$ ($r=0$: desorden; $r=1$: sincronización global).

\section*{1. Criterio operativo}
Se define $K_c$ como el menor acoplamiento para el cual el sistema queda más cerca
de la sincronización ($r=1$) que del desorden ($r=0$), es decir $r_\infty \ge \tfrac{1}{2}$:
\[
\boxed{\;K_c = \min\bigl\{\,K : r_\infty(K) \ge \tfrac{1}{2}\,\bigr\}\;}
\]
El cruce se interpola en escala $\log K$ entre los dos puntos del barrido que lo
encierran. Red completa: $K_c \approx \num{3.7e-4}$.

\section*{2. Referencia teórica (campo medio)}
La teoría de Kuramoto predice el \emph{inicio} de la sincronización en
\[
K_c^{\text{MF}} = \frac{2}{\pi\,g(0)},
\]
con $g$ la densidad de frecuencias naturales evaluada en su media. Para
$\mathcal{N}(\mu,\sigma^2)$ se tiene $g(0)=\tfrac{1}{\sigma\sqrt{2\pi}}$, y como $K$
multiplica la suma \emph{cruda} sobre vecinos (empuje efectivo $\sim K\langle k\rangle$,
con $\langle k\rangle$ el grado medio):
\[
\boxed{\;K_c^{\text{MF}} = \frac{\sigma}{\langle k\rangle}\sqrt{\frac{8}{\pi}}\;}
\]
Red completa ($\langle k\rangle = N-1 = 599$, $\sigma=\num{0.1}$):
$K_c^{\text{MF}} \approx \num{2.7e-4}$. Marca la \emph{emergencia} de la
sincronización; el $K_c$ operativo ($r_\infty=\tfrac12$) cae algo por encima, ya
dentro de la transición.

\end{document}
